\documentclass[UTF8]{ctexart}
\usepackage{amsmath, amssymb, geometry}
\usepackage{xcolor}
\usepackage{enumitem}
\usepackage{tcolorbox}
\usepackage{cases}
\usepackage{fancyhdr}
\geometry{a4paper, margin=1in}

% 定义红色环境
\newenvironment{redenv}{\color{red}}{}

% 定义详细解析环境
\newenvironment{detailedanalysis}{%
    \color{red}
    \begin{enumerate}[label=\textbf{第\arabic*步：}, leftmargin=*, align=left, itemsep=1.5ex]
}{%
    \end{enumerate}
}

% 定义概念框
\newenvironment{conceptbox}{%
    \begin{tcolorbox}[colback=red!5, colframe=red!50!black, title=概念解析, fonttitle=\bfseries]
}{%
    \end{tcolorbox}
}

% 定义公式推导环境
\newenvironment{formuladerivation}{%
    \begin{enumerate}[label={\textbf{公式\arabic*：}}, leftmargin=*, align=left]
}{%
    \end{enumerate}
}

% 定义题目和解析的分隔线
\newcommand{\problemrule}{\noindent\rule{\textwidth}{1.5pt}\vspace{-0.8em}}
\newcommand{\analysisrule}{\noindent\rule{\textwidth}{0.8pt}\vspace{-0.8em}}

% 宏定义：方便排版选择题选项
\newcommand{\choicesFour}[4]{
    \begin{enumerate}[label=\Alph*., itemsep=0pt, parsep=0pt, topsep=5pt, labelsep=0.5em]
        \begin{minipage}{0.22\linewidth} \item #1 \end{minipage}
        \begin{minipage}{0.22\linewidth} \item #2 \end{minipage}
        \begin{minipage}{0.22\linewidth} \item #3 \end{minipage}
        \begin{minipage}{0.22\linewidth} \item #4 \end{minipage}
    \end{enumerate}
}
\newcommand{\choicesTwo}[4]{
    \begin{enumerate}[label=\Alph*., itemsep=0pt, parsep=0pt, topsep=5pt, labelsep=0.5em]
        \begin{minipage}{0.45\linewidth} \item #1 \end{minipage}
        \begin{minipage}{0.45\linewidth} \item #2 \end{minipage}
        \begin{minipage}{0.45\linewidth} \item #3 \end{minipage}
        \begin{minipage}{0.45\linewidth} \item #4 \end{minipage}
    \end{enumerate}
}
\newcommand{\choices}[4]{
    \begin{enumerate}[label=\Alph*., itemsep=0pt, parsep=0pt, topsep=5pt, labelsep=0.5em]
        \item #1
        \item #2
        \item #3
        \item #4
    \end{enumerate}
}

\newcommand{\blank}[1]{\underline{\hspace{#1}}}

\begin{document}

\section*{第七章 定积分（提高）}

% ========== 第1题 ==========
\problemrule
\section*{【题目1】变上限积分求导}
$ \frac{d}{dx}[x\int_{0}^{x}\sqrt{1+t^{4}}dt]= $ \blank{1cm}。
\choicesTwo
{$ \int_{0}^{x}\sqrt{1+t^{4}}dt $}
{$ 4x^{4}\int_{0}^{x}\sqrt{1+t^{4}}dt $}
{$ \int_{0}^{x}\sqrt{1+t^{4}}dt + x\sqrt{1+x^{4}} $}
{$ x\sqrt{1+x^{4}} $}

\begin{redenv}
\analysisrule
\subsection*{逐步解析}

\begin{detailedanalysis}
\item \textbf{详细求解过程}
应用乘法法则：$(uv)' = u'v + uv'$。
令 $u = x$，$v = \int_0^x \sqrt{1+t^4} dt$。
$$ u' = 1 $$
$$ v' = (\int_0^x \sqrt{1+t^4} dt)' = \sqrt{1+x^4} $$
$$ (uv)' = 1 \cdot \int_0^x \sqrt{1+t^4} dt + x \cdot \sqrt{1+x^4} $$
即：$\int_{0}^{x}\sqrt{1+t^{4}}dt + x\sqrt{1+x^{4}}$。

\item \textbf{最终答案}
\textbf{答案}：C
\end{detailedanalysis}
\end{redenv}

% ========== 第2题 ==========
\problemrule
\section*{【题目2】变上限积分导数公式}
设 $ f(x) $ 是连续函数，则 $ \frac{d}{dx}\int_{x^{2}}^{-1}f(t)dt= $ \blank{1cm}。
\choicesFour{$ f(x^{2}) $}{$ 2xf(x^{2}) $}{$ -f(x^{2}) $}{$ -2xf(x^{2}) $}

\begin{redenv}
\analysisrule
\subsection*{逐步解析}

\begin{detailedanalysis}
\item \textbf{详细求解过程}
\begin{enumerate}[label=(\arabic*)]
\item \textbf{调整限}：
$$ \int_{x^2}^{-1} f(t) dt = - \int_{-1}^{x^2} f(t) dt $$
\item \textbf{求导}：
应用链式法则：$\frac{d}{dx} \int_a^{u(x)} f(t) dt = f(u(x)) u'(x)$。
$$ \frac{d}{dx} (- \int_{-1}^{x^2} f(t) dt) = - [f(x^2) \cdot (x^2)'] = - f(x^2) \cdot 2x = -2xf(x^2) $$
\end{enumerate}

\item \textbf{最终答案}
\textbf{答案}：D
\end{detailedanalysis}
\end{redenv}

% ========== 第3题 ==========
\problemrule
\section*{【题目3】定积分存在性}
$ f(x) $ 在 $ [a, b] $ 上连续是 $ \int_{a}^{b} f(x) \, dx $ 存在的 \blank{1cm}。
\choicesTwo
{充分必要条件}
{充分条件非必要条件}
{必要条件非充分条件}
{既非必要也非充分条件}

\begin{redenv}
\analysisrule
\subsection*{逐步解析}

\begin{detailedanalysis}
\item \textbf{详细求解过程}
\begin{enumerate}[label=(\arabic*)]
\item \textbf{充分性}：连续函数必可积。正确。
\item \textbf{必要性}：可积不一定连续。例如只有有限个间断点的有界函数是可积的。错误。
\item \textbf{结论}：连续是可积的充分非必要条件。
\end{enumerate}

\item \textbf{最终答案}
\textbf{答案}：B
\end{detailedanalysis}
\end{redenv}

% ========== 第4题 ==========
\problemrule
\section*{【题目4】变限积分含参求导}
$ \frac{d}{dx}\int_{0}^{x}g(x)f(t)dt= $ \blank{1cm}。
\choicesTwo
{$ g(x)f(x) $}
{$ g'(x)f'(x) $}
{$ g'(x)f(x) + g(x)f'(x) $}
{$ g(x)f(x) + g'(x) \int_{0}^{x} f(t) dt $}

\begin{redenv}
\analysisrule
\subsection*{逐步解析}

\begin{detailedanalysis}
\item \textbf{详细求解过程}
$g(x)$ 与积分变量 $t$ 无关，可提到积分号外。
$$ \int_0^x g(x) f(t) dt = g(x) \int_0^x f(t) dt $$
求导（乘法法则）：
$$ (g(x) \int_0^x f(t) dt)' = g'(x) \int_0^x f(t) dt + g(x) (\int_0^x f(t) dt)' $$
$$ = g'(x) \int_0^x f(t) dt + g(x) f(x) $$

\item \textbf{最终答案}
\textbf{答案}：D
\end{detailedanalysis}
\end{redenv}

% ========== 第5题 ==========
\problemrule
\section*{【题目5】定积分极值问题}
设函数 $ y = \int_{0}^{x} (t - 1) dt $ ，则 $y$ 有 \blank{1cm}。
\choicesFour{极小值 $ \frac{1}{2} $}{极小值 $ -\frac{1}{2} $}{极大值 $ \frac{1}{2} $}{极大值 $ -\frac{1}{2} $}

\begin{redenv}
\analysisrule
\subsection*{逐步解析}

\begin{detailedanalysis}
\item \textbf{详细求解过程}
\begin{enumerate}[label=(\arabic*)]
\item \textbf{求导}：$y' = x - 1$。
\item \textbf{驻点}：令 $y'=0 \implies x=1$。
\item \textbf{判定}：
$x < 1$ 时 $y' < 0$（减）；
$x > 1$ 时 $y' > 0$（增）。
所以 $x=1$ 是极小值点。
\item \textbf{计算极值}：
$$ y(1) = \int_0^1 (t-1) dt = [\frac{1}{2}t^2 - t]_0^1 = \frac{1}{2} - 1 = -\frac{1}{2} $$
\end{enumerate}

\item \textbf{最终答案}
\textbf{答案}：B
\end{detailedanalysis}
\end{redenv}

% ========== 第6题 ==========
\problemrule
\section*{【题目6】定积分换元求导}
已知 $ f(x) $ 连续，则 $ \frac{d}{dx} \int_{0}^{x} f(x - t) dt = $ \blank{2cm}。

\begin{redenv}
\analysisrule
\subsection*{逐步解析}

\begin{detailedanalysis}
\item \textbf{详细求解过程}
方法一（推荐）：先换元。
令 $u = x-t$，则 $t=x-u$，$dt = -du$。
当 $t=0, u=x$；当 $t=x, u=0$。
$$ \int_0^x f(x-t) dt = \int_x^0 f(u) (-du) = \int_0^x f(u) du $$
所以求导：
$$ \frac{d}{dx} \int_0^x f(u) du = f(x) $$

方法二：Leibniz 积分公式。
$\int_a^x f(x,t) dt$ 求导复杂，不如换元简单。

\item \textbf{最终答案}
\textbf{答案}：$ f(x) $
\end{detailedanalysis}
\end{redenv}

% ========== 第7题 ==========
\problemrule
\section*{【题目7】凑微分法计算定积分}
$ \int_{1}^{e^{3}}\frac{dx}{x\sqrt{1+\ln x}}= $ \blank{2cm}。

\begin{redenv}
\analysisrule
\subsection*{逐步解析}

\begin{detailedanalysis}
\item \textbf{详细求解过程}
令 $u = 1 + \ln x$，则 $du = \frac{1}{x} dx$。
限：$x=1 \implies u=1$；$x=e^3 \implies u=4$。
$$ \int_1^4 \frac{du}{\sqrt{u}} = \int_1^4 u^{-1/2} du = [2u^{1/2}]_1^4 = 2(2 - 1) = 2 $$

\item \textbf{最终答案}
\textbf{答案}：2
\end{detailedanalysis}
\end{redenv}

% ========== 第8题 ==========
\problemrule
\section*{【题目8】几何意义求定积分}
定积分 $ \int_{0}^{2}\sqrt{4-x^{2}}dx= $ \blank{2cm}。

\begin{redenv}
\analysisrule
\subsection*{逐步解析}

\begin{detailedanalysis}
\item \textbf{详细求解过程}
$y = \sqrt{4-x^2}$ 是圆 $x^2 + y^2 = 4$ 上半部分。
积分范围 $x \in [0, 2]$ 对应第一象限的四分之一圆。
半径 $r=2$。
面积 $S = \frac{1}{4} \pi r^2 = \frac{1}{4} \pi \cdot 4 = \pi$。

\item \textbf{最终答案}
\textbf{答案}：$ \pi $
\end{detailedanalysis}
\end{redenv}

% ========== 第9题 ==========
\problemrule
\section*{【题目9】奇偶性简化积分}
定积分 $ \int_{-1}^{1}\frac{x^{2}+x^{5}\sin x^{2}}{1+x^{2}}dx= $ \blank{2cm}。

\begin{redenv}
\analysisrule
\subsection*{逐步解析}

\begin{detailedanalysis}
\item \textbf{详细求解过程}
$$ I = \int_{-1}^1 \frac{x^2}{1+x^2} dx + \int_{-1}^1 \frac{x^5 \sin x^2}{1+x^2} dx $$
\begin{enumerate}[label=(\arabic*)]
\item \textbf{第二项}：
设 $g(x) = \frac{x^5 \sin x^2}{1+x^2}$。
$g(-x) = \frac{(-x)^5 \sin(-x)^2}{1+(-x)^2} = \frac{-x^5 \sin x^2}{1+x^2} = -g(x)$。
奇函数在对称区间积分为0。
\item \textbf{第一项}：
$\frac{x^2}{1+x^2} = 1 - \frac{1}{1+x^2}$。偶函数。
$$ I = 2 \int_0^1 (1 - \frac{1}{1+x^2}) dx = 2 [x - \arctan x]_0^1 $$
$$ = 2 (1 - \frac{\pi}{4}) = 2 - \frac{\pi}{2} $$
\end{enumerate}

\item \textbf{最终答案}
\textbf{答案}：$ 2 - \frac{\pi}{2} $
\end{detailedanalysis}
\end{redenv}

% ========== 第10题 ==========
\problemrule
\section*{【题目10】定积分公式}
定积分 $ \int_{0}^{a}x^{2}\sqrt{a^{2}-x^{2}}dx= $ \blank{2cm}。

\begin{redenv}
\analysisrule
\subsection*{逐步解析}

\begin{detailedanalysis}
\item \textbf{详细求解过程}
令 $x = a\sin t$，$dx = a\cos t dt$。
限：$t \in [0, \pi/2]$。
$$ \int_0^{\pi/2} a^2\sin^2 t \cdot a\cos t \cdot a\cos t dt = a^4 \int_0^{\pi/2} \sin^2 t \cos^2 t dt $$
利用 Wallis 公式或三角恒等变形 $\sin^2 t \cos^2 t = \frac{1}{4} \sin^2 2t = \frac{1}{8}(1-\cos 4t)$。
$$ = \frac{a^4}{8} \int_0^{\pi/2} (1-\cos 4t) dt = \frac{a^4}{8} [\frac{\pi}{2} - 0] = \frac{\pi a^4}{16} $$

\item \textbf{最终答案}
\textbf{答案}：$ \frac{\pi a^4}{16} $
\end{detailedanalysis}
\end{redenv}

% ========== 第11题 ==========
\problemrule
\section*{【题目11】函数列极限与间断点}
设函数 $ f(x)=\lim_{n\to\infty}\frac{d}{dx}\left(\int_{-1}^{x}\frac{1+t}{1+t^{4n}}dt\right) $ ，求 $ f(x) $ 的间断点。

\begin{redenv}
\analysisrule
\subsection*{逐步解析}

\begin{detailedanalysis}
\item \textbf{详细求解过程}
\begin{enumerate}[label=(\arabic*)]
\item \textbf{计算导数}：
$$ \frac{d}{dx} \int_{-1}^x \frac{1+t}{1+t^{4n}} dt = \frac{1+x}{1+x^{4n}} $$
\item \textbf{计算极限求 f(x)}：
$$ f(x) = \lim_{n \to \infty} \frac{1+x}{1+x^{4n}} $$
讨论 $x$：
\begin{itemize}
    \item $|x| < 1$: $x^{4n} \to 0$，则 $f(x) = 1+x$。
    \item $|x| > 1$: $x^{4n} \to +\infty$，分母趋于无穷，则 $f(x) = 0$。
    \item $x = 1$: $1^{4n} = 1$，则 $f(1) = \frac{2}{2} = 1$。
    \item $x = -1$: $1^{4n} = 1$，则 $f(-1) = \frac{0}{2} = 0$。
\end{itemize}
\item \textbf{分析间断点}：
\begin{itemize}
    \item $x=1$：$\lim_{x \to 1^-} (1+x) = 2$，$\lim_{x \to 1^+} 0 = 0$。跳跃间断点。
    \item $x=-1$：$\lim_{x \to -1^+} (1+x) = 0$，$\lim_{x \to -1^-} 0 = 0$。$f(-1)=0$。所以在 $x=-1$ 处连续。
\end{itemize}
\textbf{结论}：间断点为 $x=1$。
\end{enumerate}

\item \textbf{最终答案}
\textbf{答案}：$ x=1 $
\end{detailedanalysis}
\end{redenv}

% ========== 第12题 ==========
\problemrule
\section*{【题目12】含绝对值定积分}
求定积分 $ \int_{\frac{1}{e}}^{e}|\ln x|dx $。

\begin{redenv}
\analysisrule
\subsection*{逐步解析}

\begin{detailedanalysis}
\item \textbf{详细求解过程}
分段去绝对值：
$$ |\ln x| = \begin{cases} -\ln x, & \frac{1}{e} \le x < 1 \\ \ln x, & 1 \le x \le e \end{cases} $$
$$ I = \int_{1/e}^1 (-\ln x) dx + \int_1^e \ln x dx $$
原函数 $\int \ln x dx = x\ln x - x$。
第一部分：
$$ -[x\ln x - x]_{1/e}^1 = -[(0 - 1) - (\frac{1}{e}(-1) - \frac{1}{e})] = -[-1 + \frac{2}{e}] = 1 - \frac{2}{e} $$
第二部分：
$$ [x\ln x - x]_1^e = (e - e) - (0 - 1) = 1 $$
总和：
$$ 1 - \frac{2}{e} + 1 = 2 - \frac{2}{e} $$

\item \textbf{最终答案}
\textbf{答案}：$ 2 - \frac{2}{e} $
\end{detailedanalysis}
\end{redenv}

% ========== 第13题 ==========
\problemrule
\section*{【题目13】积分性质与导数}
已知 $ f(x) = e^{x^{2}} $ ，求 $ \int_{0}^{1} f'(x) f''(x) dx $。

\begin{redenv}
\analysisrule
\subsection*{逐步解析}

\begin{detailedanalysis}
\item \textbf{详细求解过程}
\begin{enumerate}[label=(\arabic*)]
\item \textbf{求不定积分}：
$$ \int f'(x) f''(x) dx = \int f'(x) d(f'(x)) = \frac{1}{2} [f'(x)]^2 $$
\item \textbf{代入限}：
$$ I = \frac{1}{2} [f'(1)]^2 - \frac{1}{2} [f'(0)]^2 $$
\item \textbf{计算 f'(x)}：
$f'(x) = e^{x^2} \cdot 2x = 2xe^{x^2}$。
$f'(1) = 2e$。
$f'(0) = 0$。
\item \textbf{计算结果}：
$$ I = \frac{1}{2} (2e)^2 - 0 = 2e^2 $$
\end{enumerate}

\item \textbf{最终答案}
\textbf{答案}：$ 2e^2 $
\end{detailedanalysis}
\end{redenv}

% ========== 第14题 ==========
\problemrule
\section*{【题目14】反三角函数积分}
求定积分 $ \int_{0}^{\frac{1}{2}}\arcsin x dx $。

\begin{redenv}
\analysisrule
\subsection*{逐步解析}

\begin{detailedanalysis}
\item \textbf{详细求解过程}
分部积分：
$$ \int \arcsin x dx = x\arcsin x - \int \frac{x}{\sqrt{1-x^2}} dx = x\arcsin x + \sqrt{1-x^2} $$
代入限 $[0, 1/2]$：
上限：$\frac{1}{2} \arcsin \frac{1}{2} + \sqrt{1 - \frac{1}{4}} = \frac{1}{2} \cdot \frac{\pi}{6} + \frac{\sqrt{3}}{2} = \frac{\pi}{12} + \frac{\sqrt{3}}{2}$。
下限：$0 + 1 = 1$。
结果：$\frac{\pi}{12} + \frac{\sqrt{3}}{2} - 1$。

\item \textbf{最终答案}
\textbf{答案}：$ \frac{\pi}{12} + \frac{\sqrt{3}}{2} - 1 $
\end{detailedanalysis}
\end{redenv}

% ========== 第15题 ==========
\problemrule
\section*{【题目15】积分不等式证明}
设 $ f(x) $ 在 $ [a,b] $ 上连续，且单调增加，证明：
$$ \int_{a}^{b}tf(t)dt\geq\frac{a+b}{2}\int_{a}^{b}f(t)dt $$

\begin{redenv}
\analysisrule
\subsection*{逐步解析}

\begin{detailedanalysis}
\item \textbf{证明过程}
利用 Chebyshev 总和不等式或构造辅助函数。
\textbf{方法一（构造法）}：
令 $C = \frac{a+b}{2}$。即证 $\int_a^b (t - C) f(t) dt \ge 0$。
$$ \int_a^b (t - \frac{a+b}{2}) f(t) dt $$
注意到 $t - \frac{a+b}{2}$ 关于区间中点 $C$ 对称（左负右正）。
但 $f(t)$ 单调增加，会导致正的部分权重大于负的部分。
严格证明：做代换 $t = a+b-u$，
$$ I = \int_b^a (a+b-u - \frac{a+b}{2}) f(a+b-u) (-du) = \int_a^b (\frac{a+b}{2} - u) f(a+b-u) du $$
即 $I = - \int_a^b (t - \frac{a+b}{2}) f(a+b-t) dt$。
所以
$$ 2I = \int_a^b (t - \frac{a+b}{2}) [f(t) - f(a+b-t)] dt $$
当 $t > \frac{a+b}{2}$ 时，$t - \frac{a+b}{2} > 0$，且 $t > a+b-t$，故 $f(t) \ge f(a+b-t)$，被积函数 $\ge 0$。
当 $t < \frac{a+b}{2}$ 时，$t - \frac{a+b}{2} < 0$，且 $t < a+b-t$，故 $f(t) \le f(a+b-t)$，被积函数 $\ge 0$。
综上，积分值 $\ge 0$。证毕。
\end{detailedanalysis}
\end{redenv}

\end{document}
